% MT59--2C
% Vergiftung mit Reizgasen
% 14.0
% 2013-11-30
\label{m:am-reizgase}
\begin{enumerate}
  \item Auf Selbstschutz achten!
  \begin{itemize}
    \item Patienten retten, aber nur wenn es der
    Eigenschutz zulässt.
    \item \ce{O2}-Berieselung mit Maske ist in keinem Fall
    ein ausreichender Atemschutz!
  \end{itemize}
  \item Gefahr erkennen: Welcher Stoff? Expertenrat einholen!
  \item Wenn erforderlich: Absperrmaßnahmen durchführen
  \item Wenn erforderlich:  Spezialkräfte anfordern
  \item Vorgehen analog zu Abschnitt \enquote{Gefahrengutunfall}
  (\prettyref{topic:gefahrengutunfall})
  \item Weitere Opfer ausschließen (lassen):
  Nachbarwohnungen, etc
  \item Sicherung der Vitalfunktionen, Beurteilung ob der
  Patient vital bedroht ist.
  \item |TxBeiVit| 
  \item Traumacheck (Folgeverletzungen)
  \item Neurocheck inkl. Blutzuckerbestimmung
  (Differentialdiagnosen, Beurteilung der Beeinträchtigung im Verlauf)
\end{enumerate}